\section{Decisions Without Generation}
\sectionkicker{DECISION MODELS}
\begin{wideflow}
{\Large\bfseries 生成しない決断：System OneとJev}\par
\smallskip
{\color{SurveyMuted}9月15日の非生成・構造化判断モデルを、仕組みと商用条件の範囲で示す。}\par
\end{wideflow}

9月15日、TypeSafe AIは新しい種類のフロンティアモデルとしてSystem Oneと、その最初の実装Jevの早期提供を発表した\cite{jev}。文章を紡ぐことを捨て、速く構造化された判断をソフトウェアが直接使える形で返す設計で、新しいモデル構成・並列サンプラー・「Reinforcement Learning for Calibrated Decisions（RLCD）」という学習法の三点が添えられる。位置づけは「フロンティア知能の関数呼び出し」で、非構造の状態を入れて型つきの確率的判断を出す。Noul・Choice・Scoreといった素子、構造化されたプログラム状態の入力、一度の問いですべての出力を並列に引く仕組み、装置を意識した作りが示される。既存の大規模言語モデル並みの知能をSystem Oneの仕事で保ちながら、2桁ほど速く効率的とされるが、独立した計測はない。

商用条件は、入力が100万トークンあたり0.042ドル、出力は計量しないほど安いため無料とされる\cite{jev}。速さは端から端で70ミリ秒から500ミリ秒、同じ知能で40倍から200倍速いとされ、業務評価では最大193.6倍速・444.6倍安いという上限側の値も示される。評価は、GPT-5.6 Terraとの並べ比べ（問いを簡単化した留保つき）や、GPT-6 AstraとFable 5.1の平均を参照点にした4業務評価（参照点の偏りの注意書きつき）などベンダー側のまとめであり、独立した再現はない。「幻覚を起こさない」は型の範囲の保証（型誤りを起こさない、確信度の目盛りが合っている）であり、誤った選択肢を選びうることまで否定しない。パラメーター数や構造、損失の作りは資料になく、作らない。公開時刻は日単位（9月15日）で押さえ、時刻は言い切らない。

\begin{claimboundary}[判断モデルの境界]
速さや費用、評価の値はベンダー側のまとめであり、独立した再現はない。型の保証を質の保証に広げない。数や構造の記述は資料にないため作らない。
\end{claimboundary}

\begin{communitynote}[X / Community observation --- context only]
Jevの発表や組み立て実演への公開投稿の反応（創業者の発表投稿、実演や経路分析の言及など）は、関心の向きを示す利用者の文脈として残す\cite{w38x-jev-ceo,w38x-jev-route}。速度や料金のまとめが投稿で回っても、構造の事実には繰り上げない。
\end{communitynote}
